\documentclass[conference]{IEEEtran}
\usepackage{cite}
\usepackage{amsmath,amssymb,amsfonts}
\usepackage{algorithmic}
\usepackage{graphicx}
\usepackage{textcomp}
\usepackage{xcolor}
\usepackage{booktabs}
\usepackage{hyperref}

\begin{document}

\title{Metal Surface Defect Classification Using Deep Learning}

\author{\IEEEauthorblockN{Dimple Parmar}
\IEEEauthorblockA{\textit{Machine Learning Course Project} \\
}}

\maketitle

\begin{abstract}
Surface defect detection is a critical process in metallurgical manufacturing to guarantee material quality and structural integrity. Traditional methods relying on manual visual inspection are inefficient, expensive, and error-prone. This study proposes an automated deep learning framework for classifying hot-rolled steel surface defects. We developed a Custom Convolutional Neural Network (CNN) trained from scratch and compared it against a transfer learning model based on the pre-trained MobileNetV2 architecture. Both models were trained and evaluated on the NEU Metal Surface Defect Dataset, containing 1,800 images distributed equally across six classes. To ensure rigorous evaluation, we implemented a stratified 70/15/15 train/validation/test split. The MobileNetV2 transfer learning model, combining locked feature extraction and fine-tuning, achieved superior classification results with an accuracy of 98.15\%, compared to 92.59\% for the Custom CNN. In addition, Gradient-weighted Class Activation Mapping (Grad-CAM) was implemented to provide spatial interpretability of the models' predictions, demonstrating that the classifier correctly isolates physical defects on the steel strips.
\end{abstract}

\section{Introduction}
In hot-rolled steel strip manufacturing, defects such as crazing, inclusions, and scratches frequently occur due to rolling roll wear, mechanical vibration, or slab impurities. The presence of these defects degrades the mechanical properties, corrosion resistance, and aesthetic appeal of the metal. Consequently, real-time quality control is vital for modern steel fabrication processes.

Early automated surface inspection (ASI) systems relied on manual visual inspection. However, human inspectors struggle to keep pace with rapid rolling speeds, and fatigue leads to inconsistent error rates. The computer vision community subsequently introduced machine learning methods combining handcrafted texture descriptors (e.g., Local Binary Patterns, Gabor filters) with Support Vector Machines (SVM). These methods, while faster, proved fragile under variations in illumination, scale shifts, and ambient noise.

In recent years, deep learning has revolutionized computer vision. Convolutional Neural Networks (CNNs) automatically learn hierarchical spatial features directly from raw pixels, minimizing the need for handcrafted descriptors. This study evaluates two distinct paradigms of deep learning for metal defect classification: (1) designing a custom sequential CNN optimized for computational efficiency, and (2) leveraging transfer learning with MobileNetV2, pre-trained on ImageNet, to classify specialized metal surfaces. We evaluate both approaches using metrics like Accuracy, Precision, Recall, F1 Score, and training compute times. We also use Grad-CAM to visualize what regions of the image influence the model's predictions.

\section{Related Work}
Automated inspection of hot-rolled steel strip defects has a rich literature. Early work focused on traditional texture analysis. Song and Yan \cite{song2013noise} proposed Completed Local Binary Patterns (CLBP) to capture defect patterns, demonstrating that texture features are susceptible to high noise levels and variations in steel surface textures.

The transition to deep learning addressed these issues. Luo et al. \cite{luo2020automated} surveyed visual defect detection methods, showing that CNN architectures outperform traditional handcrafted feature extraction in both classification accuracy and robustness against surface noise. Yi et al. \cite{yi2017end} implemented a custom sequential CNN for the NEU dataset, demonstrating that an end-to-end network could extract features and classify defects without separate preprocessing steps.

Transfer learning has since become a standard methodology. Pre-training on massive datasets like ImageNet allows models to learn general features like edges and textures, which can then be fine-tuned for specific domains, as shown by He et al. with ResNet \cite{he2016deep}. For edge computing and real-time processing, Sandler et al. proposed MobileNetV2 \cite{sandler2018mobilenetv2}, which utilizes inverted residual blocks and depthwise separable convolutions to achieve high efficiency. This research evaluates the applicability of both a custom CNN and MobileNetV2 for metal surface defect classification.

\section{Methodology}
The proposed framework consists of four primary stages: data ingestion, preprocessing/augmentation, model construction, and model interpretation.

\subsection{Dataset Description}
This study utilizes the NEU Metal Surface Defect Database. The dataset is balanced, containing 1,800 grayscale images with 300 samples for each of the six defect classes: Crazing (\texttt{crazing}), Inclusion (\texttt{inclusion}), Patches (\texttt{patches}), Pitted Surface (\texttt{pitted\_surface}), Rolled-in Scale (\texttt{rolled-in\_scale}), and Scratches (\texttt{scratches}). The original images are 200x200 pixels in BMP format.

\begin{table}[htbp]
\caption{Dataset Distribution}
\begin{center}
\begin{tabular}{lcc}
\toprule
\textbf{Defect Class} & \textbf{Number of Images} & \textbf{Percentage} \\
\midrule
Crazing & 300 & 16.67\% \\
Inclusion & 300 & 16.67\% \\
Patches & 300 & 16.67\% \\
Pitted Surface & 300 & 16.67\% \\
Rolled-in Scale & 300 & 16.67\% \\
Scratches & 300 & 16.67\% \\
\midrule
\textbf{Total} & \textbf{1,800} & \textbf{100.0\%} \\
\bottomrule
\end{tabular}
\label{tab:dataset}
\end{center}
\end{table}

\subsection{Preprocessing \& Data Splitting}
The preprocessing pipeline operates as follows:
\begin{itemize}
    \item \textbf{Directory Scan}: The raw files are loaded programmatically.
    \item \textbf{Resizing}: Images are bilinearly interpolated to 224x224 pixels to match the standard input resolution of MobileNetV2.
    \item \textbf{Normalization}: Grayscale values are normalized to [0, 1] to stabilize gradient calculations.
    \item \textbf{Stratified Splitting}: To prevent data leakage and ensure stable evaluation, the dataset is programmatically split into: Training Set (70\%), Validation Set (15\%), and Test Set (15\%).
\end{itemize}

\subsection{Data Augmentation}
To mitigate overfitting, the training generator applies real-time data augmentations, including random rotation ($\pm 20$ degrees), horizontal and vertical flips, width and height translations (10\%), random zoom ($\pm 20\%$), and brightness adjustments.

\subsection{Model Architectures}
We implement and compare two distinct architectures:
\begin{enumerate}
    \item \textbf{Custom CNN}: Constructed sequentially with four Conv2D blocks:
    \begin{itemize}
        \item Block 1-4: Convolution layers with increasing filters (32, 64, 128, 256), $3\times3$ kernels, ReLU activation, Batch Normalization, and MaxPool (2x2).
        \item Head: Flatten $\rightarrow$ Dense (512, ReLU, Dropout 0.5) $\rightarrow$ Dense (256, ReLU, Dropout 0.5) $\rightarrow$ Softmax (6).
    \end{itemize}
    \item \textbf{MobileNetV2 Transfer Learning}: Implements a pre-trained MobileNetV2 base with frozen weights. A custom classifier head is appended: Global Average Pooling $\rightarrow$ Dense (256, ReLU, Dropout 0.5) $\rightarrow$ Softmax (6). After initial training, layers from index 100 onwards are unfrozen and trained at a lower learning rate ($1\times10^{-5}$).
\end{enumerate}

\section{Experimental Setup}
\begin{itemize}
    \item \textbf{Software Stack}: Python 3.8, TensorFlow 2.16.1, Keras, OpenCV.
    \item \textbf{Loss Function}: Categorical Crossentropy.
    \item \textbf{Optimizer}: Adam (initial learning rate = 0.001).
    \item \textbf{Callbacks}: EarlyStopping, ReduceLROnPlateau, and ModelCheckpoint.
\end{itemize}

\section{Results}
Both models were trained and evaluated on the holdout test set (270 images). 

\subsection{Model Performance Summary}
The table below compares the Custom CNN and MobileNetV2 models across all evaluation metrics:

\begin{table}[htbp]
\caption{Model Performance Comparison}
\begin{center}
\begin{tabular}{lccccc}
\toprule
\textbf{Model} & \textbf{Acc.} & \textbf{Prec.} & \textbf{Recall} & \textbf{F1} & \textbf{Time} \\
\midrule
Custom CNN & 92.59\% & 0.9284 & 0.9259 & 0.9252 & 3.08 m \\
MobileNetV2 & \textbf{98.15\%} & \textbf{0.9824} & \textbf{0.9815} & \textbf{0.9815} & 5.75 m \\
\bottomrule
\end{tabular}
\label{tab:results}
\end{center}
\end{table}

\subsection{Learning Curves and Convergence}
The Custom CNN required 25 epochs to converge, exhibiting slight oscillations in validation loss. MobileNetV2 displayed rapid convergence. During the initial feature extraction phase, validation accuracy stabilized at $\sim$95.2\% within 10 epochs. After fine-tuning, test accuracy increased to 98.15\%.

\subsection{Interpretability via Grad-CAM}
Grad-CAM heatmaps were generated to localize predictions. For 'Scratches', the heatmaps correctly highlighted the linear scratch paths on the steel. For 'Patches' and 'Pitted Surface', the highest activation intensities matched the physical location of surface discoloration, verifying that the models classify defects based on relevant visual features rather than background noise.

\section{Discussion}
The results show that Transfer Learning is highly effective for specialized industrial classification tasks. Although MobileNetV2 is trained on natural images, the low-level edge and texture filters generalize well to industrial steel surfaces. The Custom CNN performed well, achieving over 92\% accuracy, but was more prone to overfitting. For resource-constrained factory environments, the Custom CNN has a smaller file size (17 MB) compared to MobileNetV2 (33 MB), but MobileNetV2's superior accuracy makes it the preferred option for quality assurance applications.

\section{Limitations and Future Work}
While the MobileNetV2 model achieved high accuracy, a key limitation of image classification is its inability to localize multiple distinct defects on a single large steel sheet. Future work will focus on transitioning from image classification to object detection architectures (e.g., YOLOv8) to draw bounding boxes around multiple defects in real-time. Deploying these models onto edge devices like the NVIDIA Jetson Nano for live factory integration is the next practical step.

\section{Conclusion}
This study developed a deep learning system for classifying hot-rolled steel surface defects using the NEU dataset. A Custom CNN trained from scratch was compared against a MobileNetV2 transfer learning model. Rigorous testing on a 15\% stratified split demonstrated that the fine-tuned MobileNetV2 achieves a test accuracy of 98.15\%, outperforming the Custom CNN (92.59\%). Grad-CAM visualizations confirmed that the models focus on the correct defect regions. 

\begin{thebibliography}{00}
\bibitem{luo2020automated} Q. Luo, X. Fang, L. Liu, C. Yang, and Y. Sun, ``Automated Visual Defect Detection for Flat Steel Surface: A Survey,'' \textit{IEEE Transactions on Instrumentation and Measurement}, vol. 69, no. 9, pp. 6261-6280, Sept. 2020.
\bibitem{song2013noise} K. Song and Y. Yan, ``A noise robust method based on completed local binary patterns for hot-rolled steel strip surface defects,'' \textit{Applied Surface Science}, vol. 285, pp. 858-864, Nov. 2013.
\bibitem{yi2017end} L. Yi, G. Li, and M. Jiang, ``End-to-end steel strip surface defects recognition based on convolutional neural networks,'' \textit{Steel Research International}, vol. 88, no. 2, p. 1600068, Feb. 2017.
\bibitem{he2016deep} K. He, X. Zhang, S. Ren, and J. Sun, ``Deep residual learning for image recognition,'' in \textit{Proceedings of the IEEE Conference on Computer Vision and Pattern Recognition (CVPR)}, 2016, pp. 770-778.
\bibitem{sandler2018mobilenetv2} M. Sandler, A. Howard, M. Zhu, A. Zhmoginov, and L. C. Chen, ``MobileNetV2: Inverted Residuals and Linear Bottlenecks,'' in \textit{Proceedings of the IEEE/CVF Conference on Computer Vision and Pattern Recognition (CVPR)}, 2018, pp. 4510-4520.
\end{thebibliography}

\end{document}
